\documentclass[11pt]{article}
\usepackage{amsmath}
\usepackage{amssymb}
\usepackage{amsthm}
\usepackage[margin=1in]{geometry}
\usepackage{hyperref}
\usepackage{listings}

\theoremstyle{definition}
\newtheorem{definition}{Definition}
\theoremstyle{plain}
\newtheorem{theorem}{Theorem}
\newtheorem{lemma}{Lemma}

\lstset{
  basicstyle=\ttfamily\small,
  breaklines=true,
  columns=flexible,
  keepspaces=true
}

\title{Conjecture 00000001559: the six-disk covering radius of the unit disk\\
\large Disproof of the stated value of $r_6$}
\author{Zhiyuan Zhang (GitHub: LibertychaserUS)}
\date{2026-10-27}

\begin{document}
\maketitle

\section{The conjecture as filed}

The statement, quoted in English from the filing (symbols typeset in math mode, wording unchanged):

\begin{quote}
Definition: The $n$-disk covering radius $r_n$ of the unit disk is the minimal radius of $n$ equal disks covering it. Conjecture: $r_6$ is an explicit quadratic algebraic number of the $(1+\sqrt{3})/2\cdot(\sqrt{3}-1)$ type, realized by one central disk and five symmetric annular disks; and $r_7$ is a cubic algebraic number.
\end{quote}

\section{Precise reading}

Let $D=\{p\in\mathbb{R}^2 : |p|\le 1\}$ be the closed unit disk and $B(c,r)=\{p : |p-c|\le r\}$ the closed disk of radius $r$ centred at $c$. The covering radius used is
\begin{equation}
  r_n \;=\; \inf\Bigl\{\, r\ge 0 \;:\; \exists\, c_1,\dots,c_n\in\mathbb{R}^2,\ D\subseteq \textstyle\bigcup_{i=1}^n B(c_i,r)\,\Bigr\}.
\end{equation}
Centres are unrestricted. Repetitions are allowed in the main reading. The infimum is attained by compactness, so it is a minimum. The disproof also holds when the $n$ centres must be pairwise distinct (see Section~\ref{sec:lean}).

The clause ``$r_6$ is of the $(1+\sqrt{3})/2\cdot(\sqrt{3}-1)$ type'' is read as the equality
\[
  r_6 \;=\; \frac{1+\sqrt{3}}{2}\,(\sqrt{3}-1).
\]
The right-hand side is exactly $1$, because $(1+\sqrt3)(\sqrt3-1)=3-1=2$. The literal claim is therefore $r_6=1$. The expression is rational, so it is not a quadratic irrational as the text describes.

\section{Counterexample and verification}

\subsection{The certificate}

Let $c_1=(\tfrac12,\tfrac12)$, $c_2=(-\tfrac12,\tfrac12)$, $c_3=(-\tfrac12,-\tfrac12)$, $c_4=(\tfrac12,-\tfrac12)$, and take six disks: these four, with $c_1$ and $c_3$ repeated. Let $p=(x,y)\in D$. Choose the corner $c=(\operatorname{sgn}(x)/2,\operatorname{sgn}(y)/2)$. Expanding gives
\begin{equation}\label{eq:identity}
  |p-c|^2 \;=\; x^2+y^2-|x|-|y|+\tfrac12 .
\end{equation}
Since $x^2+y^2\le 1$ we have $|x|\le 1$, so $x^2\le |x|$. Likewise $y^2\le |y|$. Hence $x^2+y^2-|x|-|y|\le 0$, and therefore $|p-c|^2\le \tfrac12$. Every point of $D$ is within distance $1/\sqrt2$ of one of the six centres, which gives
\[
  r_6 \;\le\; \frac{1}{\sqrt 2} \;<\; 1 .
\]
This contradicts the claimed value $r_6 = 1$.

\subsection{Computed values}

The script \texttt{check.py} (sympy 1.14.0, numpy 2.5.4) gives the following. The full output is in \texttt{check\_output.txt}.

{\small\begin{center}
\begin{tabular}{@{}lll@{}}
\hline
Part & Quantity & Value \\
\hline
A & claimed expression $(1+\sqrt3)/2\cdot(\sqrt3-1)$, exact (sympy) & $1$ \\
B & corner certificate, polar-grid maximum distance & $0.70710678 \le 1/\sqrt2$ \\
B & exact rational identity (\ref{eq:identity}), 15\,696 points in $D$ & holds; $d^2\le\tfrac12$ \\
B & largest exact $d^2$ among the samples & $3094976921/6250000000 \approx 0.4952 \le 0.5$ \\
C & hexagon of radius $0.7$ (scanner's suggestion), grid maximum & $0.7$; convex bound $0.49$ \\
D & random-restart search, coarse / fine grid (upper bound only) & $0.5556$ / $0.5585$ \\
\hline
\end{tabular}
\end{center}}

Part D is only a numerical upper bound, and it is not a proof of the value of $r_6$. It shows that $r_6$ is far below $1$, and it is consistent with the disproof. The $r_7$ clause was not checked. The refutation of the $r_6$ clause alone is enough to falsify the statement as written.

\section{Lean formalization}
\label{sec:lean}

The file is \texttt{Disproof.lean}, in namespace \texttt{TLMC1559}. It imports Mathlib. Mathlib has no covering radius, so the objects are defined directly. Distances are written in coordinates, because the product metric on $\mathbb{R}\times\mathbb{R}$ is the sup-norm and must not be used.

\begin{lstlisting}
def inDisk (p : R x R) : Prop := p.1 ^ 2 + p.2 ^ 2 <= 1

def Covers (r : R) (C : Fin 6 -> R x R) : Prop :=
  forall  p : R x R, inDisk p -> exists  i, (p.1 - (C i).1) ^ 2 + (p.2 - (C i).2) ^ 2 <= r ^ 2

def coverRadius : Set R :=
  {r | 0 <= r /\ exists  C : Fin 6 -> R x R, Covers r C}

def coverRadiusInj : Set R :=
  {r | 0 <= r /\ exists  C : Fin 6 -> R x R, Function.Injective C /\ Covers r C}

noncomputable def r6 : R := sInf coverRadius
noncomputable def r6Inj : R := sInf coverRadiusInj
\end{lstlisting}

The certificate is \texttt{C6} (the corners with two repeats) and \texttt{C6d} (the corners plus $(0,0)$ and $(3,3)$, which are pairwise distinct). The corner lemma \texttt{cover\_of\_corners} proves, by a case split on the signs of $x$ and $y$, that every $p\in D$ has a centre at squared distance at most $\tfrac12$. The bounds $|x|,|y|\le 1$ are obtained with \texttt{nlinarith}. The lemmas \texttt{covers\_C6} and \texttt{covers\_C6d} show that both certificates cover $D$ with radius $3/4$. Then \texttt{r6\_le} and \texttt{r6Inj\_le} give $r_6\le 3/4$ through \texttt{csInf\_le}, and \texttt{conjectured\_expr\_eq\_one} shows that the expression equals $1$ using $\sqrt3^2=3$ and \texttt{linear\_combination}.

The main theorem is:

\begin{lstlisting}
theorem conjecture_00000001559_fails :
    not (r6 = (1 + Real.sqrt 3) / 2 * (Real.sqrt 3 - 1)) /\
    not (r6Inj = (1 + Real.sqrt 3) / 2 * (Real.sqrt 3 - 1))
\end{lstlisting}

In the listings, ASCII stands for the Unicode symbols of the Lean source: \texttt{R} for $\mathbb{R}$, \texttt{/\textbackslash} for the conjunction, \texttt{not} for negation, \texttt{->} for the arrow.

The first conjunct is the disproof with unrestricted centres. The second is the disproof with pairwise distinct centres.

\paragraph{Build result.} Running \texttt{lake env lean} on the file gives no errors and no warnings. The line \texttt{\#print axioms} reports
\begin{center}
{\small\begin{lstlisting}
'TLMC1559.conjecture_00000001559_fails' depends on axioms: [propext, Classical.choice, Quot.sound]
\end{lstlisting}}
\end{center}
The file contains no \texttt{sorry}, \texttt{native\_decide}, \texttt{axiom}, \texttt{unsafe}, \texttt{opaque}, \texttt{implemented\_by}, \texttt{extern}, or \texttt{admit}. The distinctness of the six points of \texttt{C6d} is checked by \texttt{fin\_cases} and \texttt{norm\_num}; no \texttt{decide} is used.

\paragraph{Scope.} The formal statement refutes the stated value of $r_6$. It does not formalize the structural claims (a central disk with five annular disks) or the cubic claim for $r_7$.

\section{Reproduction}

\begin{enumerate}
  \item Verifier (about two minutes, mainly the random-restart search):
  \begin{lstlisting}
cd LibertychaserUS_submission_20261027064631
python3 check.py
  \end{lstlisting}
  \item Lean (Lean v4.33.1, Mathlib v4.33.1):
  \begin{lstlisting}
cd lean
lake exe cache get \&\& lake build
  \end{lstlisting}
  \item Report:
  \begin{lstlisting}
cd LibertychaserUS_submission_20261027064631
tectonic report.tex
  \end{lstlisting}
\end{enumerate}

\end{document}
